\begin{problem}[Perpendicular constraints]\label{ex:05:perpendicular-dimension}
\begin{enumerate}
\item If $a\ne0$ in $\R^n$, find the dimension of $\{x:a\cdot x=0\}$.
\item Find the dimension of the subspace of $\R^6$ perpendicular to $(1,1,-2,3,4,5)$ and $(0,0,1,1,0,7)$.
\item For $a\ne0$ and fixed $p\in\R^n$, find the dimension of $\{x:a\cdot x=a\cdot p\}$.
\end{enumerate}
\end{problem}
\begin{solution}
(a) The span of $a$ has dimension one, so its complement has dimension $n-1$.
(b) The two constraint vectors are independent: the first coordinate of a zero combination forces the first coefficient zero, and then the third coordinate forces the second zero. Their span has dimension two, so the complement has dimension four.
(c) Subtracting $p$ identifies the set as $p+\operatorname{span}(a)^\perp$, an affine set of dimension $n-1$.
\end{solution}

\begin{problem}[Bases of solution spaces]\label{ex:05:solution-bases}
Find a basis and dimension for each solution space:
\begin{enumerate}
\item $2x+y-z=0$, $2x+y+z=0$;
\item $x-y+z=0$;
\item $4x+7y-\pi z=0$, $2x-y+z=0$;
\item $x+y+z=0$, $x-y=0$, $y+z=0$.
\end{enumerate}
\end{problem}
\begin{solution}
(a) Subtracting gives $z=0$, then $y=-2x$. Basis $(1,-2,0)$; dimension one.
(b) Write $x=y-z$. Basis $(1,1,0),(-1,0,1)$; dimension two, with independence read from the last two coordinates.
(c) From $y=2x+z$, the first equation gives $18x+(7-\pi)z=0$. Set $z=18t$, obtaining $(x,y,z)=t(\pi-7,2\pi+4,18)$. The vector is nonzero, so it is a basis and the dimension is one.
(d) The second and third equations give $x=y$, $z=-y$; the first then gives $y=0$. The space is zero, with empty basis and dimension zero.
\end{solution}

\begin{problem}[Counting free directions]\label{ex:05:solution-dimensions}
Find the dimensions of the following homogeneous solution spaces:
\begin{enumerate}
\item $2x-3y+z=0$, $x+y-z=0$;
\item $2x+7y=0$, $x-2y+z=0$;
\item $2x-3y+z=0$, $x+y-z=0$, $3x+4y=0$, $5x+y+z=0$;
\item $x+y+z=0$, $2x+2y+2z=0$.
\end{enumerate}
\end{problem}
\begin{solution}
(a) The second equation gives $z=x+y$; the first gives $3x-2y=0$. All solutions are $t(2,3,5)$, so dimension one.
(b) Put $y=2t$; then $x=-7t$ and $z=11t$, again dimension one.
(c) The first two equations give $t(2,3,5)$ as in (a). The third gives $18t=0$, so only zero remains and the dimension is zero.
(d) There is one independent equation, and the solutions are $s(1,0,-1)+t(0,1,-1)$, so dimension two.
\end{solution}

\begin{problem}[Column spaces and perpendicular spaces]\label{ex:05:column-perp}
For $A\in\R^{m\times n}$, prove that every vector in $\operatorname{col}A$ is perpendicular to every vector in $\ker A^\top$, that $(\operatorname{col}A)^\perp=\ker A^\top$, and that their dimensions sum to $m$.
\end{problem}
\begin{solution}
If $x=Av$ and $A^\top y=0$, then $x\cdot y=v^\top A^\top y=0$. This proves one inclusion in the complement identity. Conversely, a vector perpendicular to the column space is perpendicular to each column, so every entry of $A^\top y$ is zero. The complement dimension theorem now gives $\dim\operatorname{col}A+\dim\ker A^\top=m$.
\end{solution}

\begin{problem}[Gram kernels and spaces]\label{ex:05:gram-spaces}
For a real matrix $A$, prove
\[
 \ker A^\top=\ker(AA^\top),\quad
 (\operatorname{col}A)^\perp=(\operatorname{col}(AA^\top))^\perp,\quad
 \operatorname{col}A=\operatorname{col}(AA^\top),
\]
and $\operatorname{rank}A=\operatorname{rank}(AA^\top)=\operatorname{rank}(A^\top A)$.
\end{problem}
\begin{solution}
One kernel inclusion follows by multiplying $A^\top x=0$ by $A$. For the converse, $AA^\top x=0$ gives $\norm{A^\top x}^2=x^\top AA^\top x=0$, hence $A^\top x=0$. Since $AA^\top$ is symmetric, the complement identity for column spaces gives the second equality. Taking complements again gives the third. Their dimensions are equal, and applying the same reasoning to $A^\top$ proves the remaining rank equality.
\end{solution}

\begin{problem}[A numerical Gram-rank comparison]\label{ex:05:numerical-gram}
For $A=\begin{pmatrix}1&5&6\\2&6&8\\7&1&8\end{pmatrix}$, find the ranks of $A$, $A^\top$, $AA^\top$, and $A^\top A$.
\end{problem}
\begin{solution}
Column three is the sum of columns one and two. The first two are independent because the equations $x+5y=0$, $2x+6y=0$ force $x=y=0$. Thus the rank is two. Equality under transpose and the Gram-rank theorem show that all four ranks are two. This reasoning avoids unnecessary multiplication of two larger Gram matrices while proving their exact ranks.
\end{solution}

\begin{problem}[Gram factors inside products]\label{ex:05:gram-product-ranks}
For conformable real matrices, prove
\[
 \operatorname{rank}(A^\top AB)=\operatorname{rank}(AB)
 =\operatorname{rank}(ABB^\top).
\]
Deduce $\operatorname{rank}(BA)=\operatorname{rank}(B^\top BA)$ and $\operatorname{rank}(AC)=\operatorname{rank}(ACC^\top)$.
\end{problem}
\begin{solution}
The Gram theorem and product-rank inequality give
\[
 \operatorname{rank}(AB)=\operatorname{rank}(B^\top A^\top AB)
 \leqslant\operatorname{rank}(A^\top AB)\leqslant\operatorname{rank}(AB).
\]
Likewise
\[
 \operatorname{rank}(AB)=\operatorname{rank}(ABB^\top A^\top)
 \leqslant\operatorname{rank}(ABB^\top)\leqslant\operatorname{rank}(AB).
\]
The opposite bounds force both equalities. Replace the pair $(A,B)$ by $(B,A)$ for the first deduction and by $(A,C)$ for the second. No nonzero-rank assumption is needed.
\end{solution}

\begin{problem}[Matrices in standard coordinates]\label{ex:05:standard-matrices}
Find the matrices of the following linear maps:
\begin{enumerate}
\item the projections of $\R^4$ onto its first two and its first three coordinates;
\item $x\mapsto cx$ on $\R^n$, including $c=7,-1$, and $(x,y)\mapsto(3x,3y)$;
\item $(x_1,x_2,x_3,x_4)\mapsto(x_1,x_2,0,0)$;
\item $F:\R^3\to\R^2$ with $F(e_1)=(1,-3)^\top$, $F(e_2)=(-4,2)^\top$, $F(e_3)=(3,1)^\top$;
\item $F(x_1,x_2,x_3)=(3x_1-2x_2+x_3,4x_1-x_2+5x_3)^\top$.
\end{enumerate}
\end{problem}
\begin{solution}
The matrices are, respectively,
\[
 \begin{pmatrix}1&0&0&0\\0&1&0&0\end{pmatrix},\quad
 \begin{pmatrix}1&0&0&0\\0&1&0&0\\0&0&1&0\end{pmatrix};\qquad
 cI_n\text{ and }3I_2;
\]
\[
 \operatorname{diag}(1,1,0,0);\qquad
 \begin{pmatrix}1&-4&3\\-3&2&1\end{pmatrix};\qquad
 \begin{pmatrix}3&-2&1\\4&-1&5\end{pmatrix}.
\]
Each column is the image of the matching standard vector. Multiplying by a general coordinate column verifies each displayed rule.
\end{solution}

\begin{problem}[Images in an abstract basis]\label{ex:05:abstract-matrices}
Let $(v_1,v_2,v_3)$ be a basis. Find the matrix in this basis for each specification:
\begin{enumerate}
\item $Fv_1=3v_2-v_3$, $Fv_2=v_1-2v_2+v_3$, $Fv_3=-2v_1+4v_2+5v_3$;
\item $Fv_1=3v_1$, $Fv_2=-7v_2$, $Fv_3=5v_3$;
\item $Fv_1=-2v_1+7v_3$, $Fv_2=-v_3$, $Fv_3=v_1$.
\end{enumerate}
\end{problem}
\begin{solution}
Place each image's coordinate column in the corresponding column of the matrix. The answers are
\[
 \begin{pmatrix}0&1&-2\\3&-2&4\\-1&1&5\end{pmatrix},\qquad
 \begin{pmatrix}3&0&0\\0&-7&0\\0&0&5\end{pmatrix},\qquad
 \begin{pmatrix}-2&0&1\\0&0&0\\7&-1&0\end{pmatrix}.
\]
For a general $v=\sum_j x_jv_j$, linearity makes these products the coordinate columns of $Fv$, which verifies the representation beyond its values on the basis.
\end{solution}

\begin{problem}[Different source and target bases]\label{ex:05:two-bases}
Given bases $(v_1,\ldots,v_n)$ of $V$ and $(w_1,\ldots,w_m)$ of $W$, describe and justify the matrix of a linear map $L:V\to W$ and its effect on coordinate columns.
\end{problem}
\begin{solution}
Write $L(v_j)=\sum_i a_{ij}w_i$ and set $A=(a_{ij})\in\R^{m\times n}$. If $v=\sum_j x_jv_j$, then $L(v)=\sum_i(\sum_j a_{ij}x_j)w_i$. Thus $[L(v)]_{(w_i)}=A[v]_{(v_j)}$. Taking $v=v_j$ forces column $j$ to be the displayed coordinate column, so this matrix is unique. Different domain and codomain dimensions and bases are both allowed.
\end{solution}

\begin{problem}[Diagonal actions and derivatives]\label{ex:05:diagonal-derivatives}
\begin{enumerate}
\item If a basis $(v_i)$ satisfies $L(v_i)=c_iv_i$, find the matrix of $L$.
\item Find the derivative matrix on $\operatorname{span}(\cos t,\sin t)$ in the ordered basis $(\cos t,\sin t)$.
\item Find the derivative matrix on $P_2$ in $(1,t,t^2)$.
\end{enumerate}
\end{problem}
\begin{solution}
(a) Column $i$ is $c_ie_i$, so the matrix is $\operatorname{diag}(c_1,\ldots,c_n)$.
(b) The derivatives are $-\sin t$ and $\cos t$, giving $\begin{pmatrix}0&1\\-1&0\end{pmatrix}$.
(c) The derivatives are $0,1,2t$, giving $\begin{pmatrix}0&1&0\\0&0&2\\0&0&0\end{pmatrix}$. The order of the functions specifies the order of coordinate entries.
\end{solution}

\begin{problem}[Polynomial identities for mappings]\label{ex:05:map-polynomials}
Let $F,L:V\to V$ be linear and satisfy $FL=LF$. Prove $(F+L)^2=F^2+2FL+L^2$, $(F-L)^2=F^2-2FL+L^2$, and $(F+L)(F-L)=F^2-L^2$. For any linear $F$ and integer $r\geqslant0$, prove $(I-F)(I+F+\cdots+F^r)=I-F^{r+1}$.
\end{problem}
\begin{solution}
Distributing the first product gives $F^2+FL+LF+L^2$; commuting combines the middle terms. Changing the signs gives the second identity, while $F^2-FL+LF-L^2=F^2-L^2$ gives the third. For the last identity, subtract $F+F^2+\cdots+F^{r+1}$ from $I+F+\cdots+F^r$; all intermediate powers cancel. The same calculation works in the reverse order because powers of a single mapping commute.
\end{solution}

\begin{problem}[Complementary projections]\label{ex:05:projection-algebra}
\begin{enumerate}
\item For linear $T$ with $T^2=I$, set $P=(I+T)/2$, $Q=(I-T)/2$. Prove $P^2=P$, $Q^2=Q$, $P+Q=I$, $\ker P=\operatorname{im}Q$ and $\operatorname{im}P=\ker Q$.
\item More generally, if $P^2=P$, set $Q=I-P$ and prove the same assertions involving $P,Q$.
\item Prove $V=\ker P\oplus\operatorname{im}P$ and uniqueness of the two summands. Verify that $(u,w)\mapsto(u,0)$ on $U\times W$ is a projection.
\end{enumerate}
\end{problem}
\begin{solution}
(a) Expanding $(I\pm T)^2/4$ and using $T^2=I$ gives $(I\pm T)/2$. Their sum is $I$, and their product in either order is $(I-T^2)/4=0$.
(b) From $P^2=P$, one has $Q^2=I-2P+P^2=Q$, $PQ=QP=0$ and $P+Q=I$. If $v=Px$, then $Qv=0$. Conversely $Qv=0$ gives $v=Pv$, hence $v\in\operatorname{im}P$. Interchanging $P,Q$ gives the other kernel-image equality, which also proves the remaining assertions in (a).
(c) Every $v=Pv+Qv$ has the required components. If $v=Px$ and $Pv=0$, then $v=P^2x=Pv=0$, so the intersection is zero. Differences of two decompositions lie in that intersection, proving uniqueness. Finally the product-space map is linear componentwise and sends $(u,0)$ to itself, so its square equals itself. Its codomain is $U\times W$, which makes its square defined.
\end{solution}

\begin{problem}[Rank of a composition]\label{ex:05:composition-rank}
For linear maps $L:V\to W$ and $F:W\to U$ of finite-dimensional spaces, prove
$\dim\operatorname{im}(FL)\leqslant\min(\dim\operatorname{im}F,\dim\operatorname{im}L)$. Deduce the matrix product-rank inequalities.
\end{problem}
\begin{solution}
The image of $FL$ is contained in the image of $F$, proving one bound. A basis $w_1,\ldots,w_r$ of $\operatorname{im}L$ has images $F(w_i)$ spanning $F(\operatorname{im}L)=\operatorname{im}(FL)$, proving dimension at most $r$. Matrix multiplication represents composition, and its image dimension is its column rank, giving both matrix inequalities.
\end{solution}

\begin{problem}[Translation and a matrix map]\label{ex:05:translation-order}
Let $A\in\R^{n\times n}$ and $c\in\R^n$, with $L_A(x)=Ax$ and $T_c(x)=x+c$. Write formulas for $L_AT_c$ and $T_cL_A$, and exhibit a case where they differ.
\end{problem}
\begin{solution}
They are $L_AT_c(x)=Ax+Ac$ and $T_cL_A(x)=Ax+c$. For example, take $A=2I_n$ and $c\ne0$; the constant terms are $2c$ and $c$, so the mappings differ. In general they coincide exactly when $Ac=c$.
\end{solution}

\begin{problem}[Inverse laws]\label{ex:05:inverse-laws}
Find the inverse of counterclockwise rotation through $\theta$ and its matrix. For invertible linear maps $F,G:V\to V$, prove $(FG)^{-1}=G^{-1}F^{-1}$.
\end{problem}
\begin{solution}
The inverse rotation is through $-\theta$, with matrix
$\begin{pmatrix}\cos\theta&\sin\theta\\-\sin\theta&\cos\theta\end{pmatrix}$.
Multiplication in either order with the rotation matrix gives $I_2$ using $\cos^2\theta+\sin^2\theta=1$.
For the mapping identity,
$(FG)(G^{-1}F^{-1})=F(GG^{-1})F^{-1}=I$ and
$(G^{-1}F^{-1})(FG)=G^{-1}(F^{-1}F)G=I$.
Thus the proposed map is the two-sided inverse, and uniqueness gives the result.
\end{solution}

\begin{problem}[Inverting coordinate formulas]\label{ex:05:inverse-formulas}
Prove invertibility and find an inverse for each map:
\begin{enumerate}
\item $(x,y)\mapsto(x+y,x-y)$;
\item $(x,y)\mapsto(2x+y,3x-5y)$;
\item $(x,y,z)\mapsto(x-y,x+z,x+y+3z)$;
\item $(x,y,z)\mapsto(2x-y+z,x+y,3x+y+z)$.
\end{enumerate}
\end{problem}
\begin{solution}
Given output $(u,v)$, (a) has inverse $((u+v)/2,(u-v)/2)$.
For (b), $y=u-2x$ and $v=3x-5(u-2x)$ give $x=(5u+v)/13$, $y=(3u-2v)/13$.
For (c), $y=x-u$, $z=v-x$, and $w=x+y+3z=3v-u-x$, so
\[
 (x,y,z)=(3v-u-w,\ 3v-2u-w,\ u-2v+w).
\]
For (d), subtracting the first output from the third gives $w-u=x+2y$, while $v=x+y$. Thus $y=w-u-v$, $x=u+2v-w$, and the first equation yields $z=-2u-5v+3w$.
Each derivation gives exactly one input for every output; substitution in the original formulas verifies existence, so each map is bijective and has the displayed inverse.
\end{solution}

\begin{problem}[Inverses from a vanishing power]\label{ex:05:nilpotent-inverses}
Let $L:V\to V$ be linear.
\begin{enumerate}
\item If $L^r=0$ for a positive integer $r$, prove $I-L$ is invertible; include the cases $r=2,3$ explicitly.
\item If $L^2+2L+I=0$, prove $L$ is invertible.
\end{enumerate}
\end{problem}
\begin{solution}
(a) The finite geometric identity in both orders gives
$(I-L)(I+L+\cdots+L^{r-1})=(I+L+\cdots+L^{r-1})(I-L)=I-L^r=I$.
For $r=2$ the inverse is $I+L$, and for $r=3$ it is $I+L+L^2$. This does not require finite-dimensional $V$.
(b) Rearrange the hypothesis as $L(-L-2I)=(-L-2I)L=I$. Thus $L^{-1}=-L-2I$, again by a two-sided identity.
\end{solution}

\begin{problem}[A basis adapted to a square-zero map]\label{ex:05:square-zero-basis}
Let $\dim V=2$, $L^2=0$, and $L\ne0$. Choose $v$ with $L(v)\ne0$ and put $w=L(v)$. Prove that $(v,w)$ is a basis of $V$.
\end{problem}
\begin{solution}
If $av+bw=0$, apply $L$: $aw+bL^2v=aw=0$. Since $w\ne0$, $a=0$, and then $bw=0$ gives $b=0$. Thus $v,w$ are independent, and their number equals $\dim V$, so they form a basis.
\end{solution}

\begin{problem}[A direct sum as a product]\label{ex:05:sum-product}
If $V=U\oplus W$, prove that $L:U\times W\to V$, $L(u,w)=u+w$, is a bijective linear mapping.
\end{problem}
\begin{solution}
Componentwise operations give linearity. The equality $V=U+W$ gives surjectivity. If $u+w=0$, then $u=-w\in U\cap W=\{0\}$, so both components are zero; the kernel is zero and $L$ is injective. Its inverse assigns the unique two components of each vector.
\end{solution}

\subsection*{Block matrix applications}
All blocks below are real and have compatible sizes. For $Z=\begin{pmatrix}A&B\\C&D\end{pmatrix}$, the default sizes are $A:m\times p$, $B:m\times q$, $C:n\times p$, $D:n\times q$; additional squareness or invertibility hypotheses are stated when required.

\begin{problem}[Block arithmetic and shape]\label{ex:05:block-arithmetic}
\begin{enumerate}
\item State the size conditions for adding two $2\times2$ block matrices, and prove that their sum is obtained by adding matching blocks.
\item Prove the block multiplication rule, specifying the matching inner block sizes.
\item Prove $Z^\top=\begin{pmatrix}A^\top&C^\top\\B^\top&D^\top\end{pmatrix}$. If $A,D$ are square, prove $\tr Z=\tr A+\tr D$.
\item When $A,D$ are square, show that symmetry, diagonality, or upper triangularity of $Z$ implies the same property for $A,D$. Find the additional conditions for each converse.
\end{enumerate}
\end{problem}
\begin{solution}
(a) Corresponding blocks must have equal sizes. Addition is entrywise, so grouping the resulting entries gives the sum blockwise.
(b) If the first matrix has row block sizes $m,n$ and column sizes $p,q$, the second must have row sizes $p,q$ and may have column sizes $r,s$. Split each dot-product sum into the first $p$ indices and the final $q$ indices. The four resulting formulas are $AE+BG$, $AF+BH$, $CE+DG$, $CF+DH$.
(c) Transpose interchanges the row and column index, putting the transposed off-diagonal blocks in the opposite positions. When diagonal blocks are square, the diagonal entries of $Z$ are exactly those of $A$ and $D$, giving the trace formula.
(d) Restricting the relevant entry equalities or zero pattern to a diagonal block proves necessity. Conversely, symmetric $A,D$ give symmetric $Z$ exactly when $C=B^\top$; diagonal $A,D$ give diagonal $Z$ exactly when $B=C=0$; upper triangular $A,D$ give upper triangular $Z$ exactly when $C=0$.
\end{solution}

\begin{problem}[Elementary operations on whole blocks]\label{ex:05:block-operations}
Verify the following left and right multiplication rules, with the identity sizes fixed by the default partition:
\[
 \begin{pmatrix}0&I_n\\I_m&0\end{pmatrix}Z=\begin{pmatrix}C&D\\A&B\end{pmatrix},\quad
 \begin{pmatrix}E&0\\0&I_n\end{pmatrix}Z=\begin{pmatrix}EA&EB\\C&D\end{pmatrix},
\]
\[
 \begin{pmatrix}I_m&E\\0&I_n\end{pmatrix}Z=\begin{pmatrix}A+EC&B+ED\\C&D\end{pmatrix},
\]
\[
 Z\begin{pmatrix}0&I_p\\I_q&0\end{pmatrix}=\begin{pmatrix}B&A\\D&C\end{pmatrix},\quad
 Z\begin{pmatrix}F&0\\0&I_q\end{pmatrix}=\begin{pmatrix}AF&B\\CF&D\end{pmatrix},
\]
\[
 Z\begin{pmatrix}I_p&F\\0&I_q\end{pmatrix}=\begin{pmatrix}A&B+AF\\C&D+CF\end{pmatrix}.
\]
\end{problem}
\begin{solution}
Apply the four block product formulas. In the first rule, for example, the upper block row is $0(A,B)+I_n(C,D)=(C,D)$, and the lower is $(A,B)$. In the second, the upper row is $(EA,EB)$ and the lower unchanged. The third adds $E$ times the lower row to the upper. The right-multiplication rules analogously interchange, scale, or add block columns. The scaling $E$ is $m\times m$ and the row-addition $E$ is $m\times n$; the scaling $F$ is $p\times p$ and the column-addition $F$ is $p\times q$. The repeated letters denote separate conformable choices, not one matrix with conflicting sizes.
\end{solution}

\begin{problem}[Block involutions and commuting shears]\label{ex:05:block-involutions}
\begin{enumerate}
\item Show that $Z=\begin{pmatrix}-I_m&B\\C&I_n\end{pmatrix}$ satisfies $Z^2=I$ if $B=0$ or $C=0$. Is this condition necessary?
\item Prove that $\begin{pmatrix}I_m&B_1\\0&I_n\end{pmatrix}$ and $\begin{pmatrix}I_m&B_2\\0&I_n\end{pmatrix}$ commute.
\end{enumerate}
\end{problem}
\begin{solution}
(a) Squaring gives $\operatorname{diag}(I_m+BC,I_n+CB)$, so the exact condition is $BC=CB=0$. The proposed condition is sufficient but unnecessary. For $m=n=2$, take $B=C=\begin{pmatrix}0&1\\0&0\end{pmatrix}$; both are nonzero and both products vanish.
(b) Either product equals $\begin{pmatrix}I_m&B_1+B_2\\0&I_n\end{pmatrix}$, since addition of the off-diagonal blocks is commutative.
\end{solution}

\begin{problem}[The two Schur eliminations]\label{ex:05:schur-eliminations}
\begin{enumerate}
\item Assuming $A$ is square and invertible, verify separately the left and right operations that clear $C$ and $B$, and their combined result $\operatorname{diag}(A,D-CA^{-1}B)$.
\item Assuming instead that $D$ is square and invertible, verify the analogous operations and their combined result $\operatorname{diag}(A-BD^{-1}C,D)$.
\end{enumerate}
\end{problem}
\begin{solution}
(a) Left multiplication by $L=\begin{pmatrix}I_m&0\\-CA^{-1}&I_n\end{pmatrix}$ gives $LZ=\begin{pmatrix}A&B\\0&D-CA^{-1}B\end{pmatrix}$. Right multiplication by $R=\begin{pmatrix}I_m&-A^{-1}B\\0&I_q\end{pmatrix}$ gives $ZR=\begin{pmatrix}A&0\\C&D-CA^{-1}B\end{pmatrix}$. Multiplying $LZ$ by $R$ gives the claimed diagonal block matrix.
(b) Set $L=\begin{pmatrix}I_m&-BD^{-1}\\0&I_n\end{pmatrix}$ and $R=\begin{pmatrix}I_p&0\\-D^{-1}C&I_n\end{pmatrix}$. The products are
$LZ=\begin{pmatrix}A-BD^{-1}C&0\\C&D\end{pmatrix}$ and
$ZR=\begin{pmatrix}A-BD^{-1}C&B\\0&D\end{pmatrix}$, whence the combined product is $\operatorname{diag}(A-BD^{-1}C,D)$. Each multiplier is invertible by changing the sign of its off-diagonal block.
\end{solution}

\begin{problem}[Inverses with two zero blocks]\label{ex:05:two-zero-inverse}
\begin{enumerate}
\item Find the inverse of $\operatorname{diag}(A,D)$ when $A,D$ are invertible.
\item Find the inverse of $Z=\begin{pmatrix}0&B\\C&0\end{pmatrix}$ when $B,C$ are square and invertible. Show that $\begin{pmatrix}0&B\\B^{-1}&0\end{pmatrix}$ is its own inverse.
\item A real square matrix is called orthogonal when its transpose is its inverse. Determine when each matrix in (b) is orthogonal.
\end{enumerate}
\end{problem}
\begin{solution}
(a) The inverse is $\operatorname{diag}(A^{-1},D^{-1})$, because either product is block diagonal with identity blocks.
(b) The inverse is $\begin{pmatrix}0&C^{-1}\\B^{-1}&0\end{pmatrix}$: multiplying in either order gives the appropriate identity diagonal blocks. Substituting $C=B^{-1}$ gives the second assertion. The off-diagonal square blocks may have different sizes; the inverse reverses the row/column partition accordingly.
(c) Comparing the transpose with the inverse gives $C^\top=C^{-1}$ and $B^\top=B^{-1}$, so both $B,C$ must be orthogonal, and this is sufficient. In the special case $C=B^{-1}$, this holds exactly when $B$ is orthogonal.
\end{solution}

\begin{problem}[Triangular block inverses]\label{ex:05:triangular-inverse}
Find the inverses of $\begin{pmatrix}I&B\\0&I\end{pmatrix}$ and $\begin{pmatrix}I&0\\C&I\end{pmatrix}$. More generally, for square invertible $A,D$, find inverses of $\begin{pmatrix}A&B\\0&D\end{pmatrix}$ and $\begin{pmatrix}A&0\\C&D\end{pmatrix}$.
\end{problem}
\begin{solution}
For identity diagonal blocks, change the sign of the single off-diagonal block. More generally the answers are
\[
 \begin{pmatrix}A^{-1}&-A^{-1}BD^{-1}\\0&D^{-1}\end{pmatrix},\qquad
 \begin{pmatrix}A^{-1}&0\\-D^{-1}CA^{-1}&D^{-1}\end{pmatrix}.
\]
For the first, the upper-right block of the product in the original order is $-BD^{-1}+BD^{-1}=0$; in the reverse order it is $A^{-1}B-A^{-1}B=0$. All diagonal products are identities and the lower-left block is zero. The second verification gives lower-left blocks $CA^{-1}-CA^{-1}=0$ and $-D^{-1}C+D^{-1}C=0$, with identity diagonal blocks.
\end{solution}

\begin{problem}[A zero diagonal block]\label{ex:05:zero-diagonal-inverse}
For invertible square $B,C$, find inverses of
\[
 Z_1=\begin{pmatrix}A&B\\C&0\end{pmatrix},\qquad
 Z_2=\begin{pmatrix}0&B\\C&D\end{pmatrix}.
\]
Include the special cases where $B,C$ are identity matrices of their respective sizes.
\end{problem}
\begin{solution}
The answers are
\[
 Z_1^{-1}=\begin{pmatrix}0&C^{-1}\\B^{-1}&-B^{-1}AC^{-1}\end{pmatrix},\quad
 Z_2^{-1}=\begin{pmatrix}-C^{-1}DB^{-1}&C^{-1}\\B^{-1}&0\end{pmatrix}.
\]
For $Z_1$, multiplication in the forward order gives diagonal blocks $BB^{-1}$ and $CC^{-1}$, upper-right block $AC^{-1}-AC^{-1}=0$, and zero lower-left. Reverse multiplication gives diagonal identities and lower-left block $B^{-1}A-B^{-1}A=0$. For $Z_2$, the corresponding cancellation is $-DB^{-1}+DB^{-1}=0$ in the forward lower-left block and $-C^{-1}D+C^{-1}D=0$ in the reverse upper-right block. Setting $B=I_m,C=I_n$ gives the two identity-block formulas, with inverse partition reversed if $m\ne n$.
\end{solution}

\begin{problem}[A scalar border]\label{ex:05:scalar-border}
\begin{enumerate}
\item Let $A$ be invertible and $\varepsilon=\delta-c^\top A^{-1}b\ne0$. Prove
\[
 \begin{pmatrix}A&b\\c^\top&\delta\end{pmatrix}^{-1}
 =\begin{pmatrix}A^{-1}&0\\0&0\end{pmatrix}
 +\frac1\varepsilon\begin{pmatrix}A^{-1}b\\-1\end{pmatrix}(c^\top A^{-1},-1).
\]
\item If $D$ is invertible and $\phi=\alpha-b^\top D^{-1}c\ne0$, prove
\[
 \begin{pmatrix}\alpha&b^\top\\c&D\end{pmatrix}^{-1}
 =\begin{pmatrix}0&0\\0&D^{-1}\end{pmatrix}
 +\frac1\phi\begin{pmatrix}-1\\D^{-1}c\end{pmatrix}(-1,b^\top D^{-1}).
\]
\end{enumerate}
\end{problem}
\begin{solution}
(a) Eliminating the lower-left block gives Schur complement $\varepsilon$. In the four-block inverse formula, replace $B,C,D,S$ by $b,c^\top,\delta,\varepsilon$. The resulting blocks are
$A^{-1}+A^{-1}bc^\top A^{-1}/\varepsilon$, $-A^{-1}b/\varepsilon$, $-c^\top A^{-1}/\varepsilon$, and $1/\varepsilon$, exactly the displayed outer-product expression.
(b) Clear the upper-right block using $D^{-1}$, obtaining the diagonal blocks $\phi,D$. Reversing the two triangular elimination factors yields blocks
$1/\phi$, $-b^\top D^{-1}/\phi$, $-D^{-1}c/\phi$, and $D^{-1}+D^{-1}cb^\top D^{-1}/\phi$. Expanding the asserted outer product gives those same four blocks. The hypotheses make every inverse used in these factorizations defined.
\end{solution}

\begin{problem}[The two block inverse formulas]\label{ex:05:schur-inverses}
For square $Z=\begin{pmatrix}A&B\\C&D\end{pmatrix}$, put $S=D-CA^{-1}B$ when $A$ is invertible, and $T=A-BD^{-1}C$ when $D$ is invertible.
\begin{enumerate}
\item Assuming $A,S$ invertible, derive the inverse using $S$.
\item Assuming $D,T$ invertible, derive the inverse using $T$.
\end{enumerate}
\end{problem}
\begin{solution}
(a) Block elimination gives $LZR=\operatorname{diag}(A,S)$, hence
\[ Z^{-1}=R\operatorname{diag}(A^{-1},S^{-1})L. \]
Multiplying yields
\[
 Z^{-1}=\begin{pmatrix}
 A^{-1}+A^{-1}BS^{-1}CA^{-1}&-A^{-1}BS^{-1}\\
 -S^{-1}CA^{-1}&S^{-1}
 \end{pmatrix}.
\]
(b) Now use $L=\begin{pmatrix}I&-BD^{-1}\\0&I\end{pmatrix}$ and $R=\begin{pmatrix}I&0\\-D^{-1}C&I\end{pmatrix}$, so $LZR=\operatorname{diag}(T,D)$. Reversing the factorization and multiplying gives
\[
 Z^{-1}=\begin{pmatrix}
 T^{-1}&-T^{-1}BD^{-1}\\
 -D^{-1}CT^{-1}&D^{-1}+D^{-1}CT^{-1}BD^{-1}
 \end{pmatrix}.
\]
Each expression is obtained by inverting a product of invertible factors, which verifies the inverse in both orders.
\end{solution}

\begin{problem}[Comparing two inverse descriptions]\label{ex:05:inverse-identity}
When all displayed inverses exist, prove
\[
 (A-BD^{-1}C)^{-1}=A^{-1}+A^{-1}B(D-CA^{-1}B)^{-1}CA^{-1},
\]
\[
 (D-CA^{-1}B)^{-1}=D^{-1}+D^{-1}C(A-BD^{-1}C)^{-1}BD^{-1}.
\]
\end{problem}
\begin{solution}
Both block inverse formulas in the preceding exercise apply to the same matrix $Z$. The inverse is unique, so its upper-left blocks must be equal, giving the first identity. Equating its lower-right blocks gives the second. The assumption that all displayed inverses exist includes both diagonal blocks and both Schur complements.
\end{solution}

\begin{problem}[A Gram matrix with two column groups]\label{ex:05:gram-block-inverse}
Let $(A\ B)\in\R^{n\times(k+m)}$ have full column rank. Define
\[
 X=A^\top A,\quad Y=A^\top B,\quad Z=B^\top B,\quad
 G=\begin{pmatrix}X&Y\\Y^\top&Z\end{pmatrix},
\]
\[
 M_A=I-AX^{-1}A^\top,\quad M_B=I-BZ^{-1}B^\top,
 \qquad E=B^\top M_AB,\quad F=A^\top M_BA.
\]
Express $G^{-1}$ first using $E^{-1}$ and then using $F^{-1}$, and justify existence of the inverses.
\end{problem}
\begin{solution}
Full column rank of $(A\ B)$ implies full column rank of each sublist. Thus $X,Z,G$ are invertible by the Gram-kernel identity and the square zero-kernel criterion. Since $E=Z-Y^\top X^{-1}Y$ and $F=X-YZ^{-1}Y^\top$, the block-elimination rank criterion implies both $E$ and $F$ are invertible. Applying the two formulas gives
\[
 G^{-1}=\begin{pmatrix}
 X^{-1}+X^{-1}YE^{-1}Y^\top X^{-1}&-X^{-1}YE^{-1}\\
 -E^{-1}Y^\top X^{-1}&E^{-1}
 \end{pmatrix},
\]
\[
 G^{-1}=\begin{pmatrix}
 F^{-1}&-F^{-1}YZ^{-1}\\
 -Z^{-1}Y^\top F^{-1}&Z^{-1}+Z^{-1}Y^\top F^{-1}YZ^{-1}
 \end{pmatrix}.
\]
Substituting the definitions of $X,Y,Z$ gives the formulas entirely in terms of $A,B,E,F$. This argument uses the given full column rank; an arbitrary pair of column groups would not justify the inverses.
\end{solution}

\begin{problem}[Three block rows]\label{ex:05:three-block}
Let
\[
 Z=\begin{pmatrix}A&B&C\\B^\top&D&0\\C^\top&0&E\end{pmatrix}
\]
be symmetric, with $D,E$ invertible, and set $Q=A-BD^{-1}B^\top-CE^{-1}C^\top$.
\begin{enumerate}
\item If $Q$ is invertible, find $Z^{-1}$.
\item Without requiring $Q$ invertible, prove $\operatorname{rank}Z=\operatorname{rank}D+\operatorname{rank}E+\operatorname{rank}Q$.
\end{enumerate}
\end{problem}
\begin{solution}
Group the last two block rows and columns. The lower-right grouped block is $H=\operatorname{diag}(D,E)$ with inverse $\operatorname{diag}(D^{-1},E^{-1})$, and its Schur complement is $Q$. This proves the rank formula by block elimination and rank additivity for block diagonal matrices.
For the inverse, set $U=D^{-1}B^\top$ and $V=E^{-1}C^\top$. Since $D,E$ are symmetric, their inverses are symmetric. The second block inverse formula gives
\[
 Z^{-1}=\begin{pmatrix}
 Q^{-1}&-Q^{-1}U^\top&-Q^{-1}V^\top\\
 -UQ^{-1}&D^{-1}+UQ^{-1}U^\top&UQ^{-1}V^\top\\
 -VQ^{-1}&VQ^{-1}U^\top&E^{-1}+VQ^{-1}V^\top
 \end{pmatrix}.
\]
Expanding $U,V$ recovers every requested block. The grouped factorization proves the identity for all nine blocks simultaneously and verifies a two-sided inverse.
\end{solution}

\begin{problem}[A bordered rectangular matrix]\label{ex:05:bordered-rank-inverse}
Let $A\in\R^{m\times n}$, $a\in\R^n$, and $\alpha\in\R$, and set $Z=\begin{pmatrix}0&A\\\alpha&a^\top\end{pmatrix}$, where the upper-left zero is an $m$-component column.
\begin{enumerate}
\item Prove $\operatorname{rank}Z=\operatorname{rank}A$ if $\alpha=0$ and $a\in\operatorname{col}A^\top$, and $\operatorname{rank}Z=\operatorname{rank}A+1$ otherwise.
\item If $m=n$, $A$ is invertible and $\alpha\ne0$, find $Z^{-1}$.
\end{enumerate}
\end{problem}
\begin{solution}
(a) If $\alpha=0$, the zero first column can be discarded. The remaining matrix appends row $a^\top$ to $A$. Its row-space dimension is unchanged if that row lies in the old row space, and increases by one otherwise. If $\alpha\ne0$, subtract $(a_j/\alpha)$ times the first column from column $j+1$ for each $j$. The last row becomes $(\alpha,0,\ldots,0)$, so after a column-block swap the matrix is block diagonal with $A$ and $\alpha$, giving rank $\operatorname{rank}A+1$.
(b) The inverse is
\[
 Z^{-1}=\begin{pmatrix}-a^\top A^{-1}/\alpha&1/\alpha\\A^{-1}&0\end{pmatrix}.
\]
In the forward product the upper-left block is $AA^{-1}=I$, upper-right is zero, lower-left is $-a^\top A^{-1}+a^\top A^{-1}=0$, and lower-right is one. The reverse product has scalar upper-left one, upper-right $-a^\top/\alpha+a^\top/\alpha=0$, zero lower-left, and lower-right $A^{-1}A=I$.
\end{solution}

\begin{problem}[Powers of a triangular block matrix]\label{ex:05:block-powers}
Let $A,D$ be square and $Z=\begin{pmatrix}A&B\\0&D\end{pmatrix}$.
\begin{enumerate}
\item For positive integer $k$, prove $Z^k=\begin{pmatrix}A^k&Q_k\\0&D^k\end{pmatrix}$ where $Q_k=\sum_{j=1}^k A^{k-j}BD^{j-1}$.
\item Simplify $Q_k$ when $D=I$ and $I-A$ is invertible.
\item Determine the corresponding expression for negative powers when they exist.
\end{enumerate}
\end{problem}
\begin{solution}
(a) The case $k=1$ is the original matrix. Multiplying $Z$ by the expression for $Z^k$ gives off-diagonal block $AQ_k+BD^k=\sum_{j=1}^{k+1}A^{k+1-j}BD^{j-1}=Q_{k+1}$ and the required diagonal powers. This proves the formula by induction.
(b) Now $Q_k=(I+A+\cdots+A^{k-1})B=(I-A)^{-1}(I-A^k)B$ by the finite geometric identity.
(c) The matrix is invertible exactly when both $A,D$ are invertible, by the block triangular rank criterion or its explicit inverse. Applying (a) to $Z^{-1}$ gives, for $k>0$,
\[
 Z^{-k}=\begin{pmatrix}A^{-k}&R_k\\0&D^{-k}\end{pmatrix},\qquad
 R_k=-\sum_{j=1}^k A^{-(k-j+1)}BD^{-j}.
\]
The positive-index sum formula cannot simply be read with a negative upper limit.
\end{solution}

\begin{problem}[Zero blocks and rank]\label{ex:05:zero-block-ranks}
\begin{enumerate}
\item Prove $\operatorname{rank}\begin{pmatrix}A&0\\0&D\end{pmatrix}=\operatorname{rank}A+\operatorname{rank}D$ and the analogous formula for $\begin{pmatrix}0&B\\C&0\end{pmatrix}$, without requiring square blocks.
\item Give counterexamples to the same diagonal-rank formula for $Z_1=\begin{pmatrix}A&B\\0&D\end{pmatrix}$ and $Z_2=\begin{pmatrix}A&0\\C&D\end{pmatrix}$.
\item Prove the formula for both $Z_1,Z_2$ if either $A$ or $D$ is square and invertible. Is that condition necessary?
\end{enumerate}
\end{problem}
\begin{solution}
(a) The column space of the diagonal block matrix is $\operatorname{col}A\times\operatorname{col}D$, with dimension the sum of component dimensions. Swapping column groups gives the off-diagonal formula.
(b) Set the scalar blocks $A=D=0$ and take $B=1$ for $Z_1$ or $C=1$ for $Z_2$. The whole matrix has rank one while both diagonal ranks vanish.
(c) If $A$ is invertible, clear $B$ by column addition in $Z_1$, or clear $C$ by row addition in $Z_2$. If $D$ is invertible, clear $B$ by row addition or clear $C$ by column addition. Each reduces to the diagonal block matrix using invertible operations, proving the equality. The condition is not necessary: if $B=C=0$, the equality holds even when both diagonal blocks are singular.
\end{solution}

\begin{problem}[Lower bounds from a zero block]\label{ex:05:zero-block-bounds}
Prove
\[
 \operatorname{rank}\begin{pmatrix}A&B\\0&D\end{pmatrix}\geqslant\operatorname{rank}A+\operatorname{rank}D,\qquad
 \operatorname{rank}\begin{pmatrix}A&0\\C&D\end{pmatrix}\geqslant\operatorname{rank}A+\operatorname{rank}D.
\]
Show that the inequality need not hold if all four blocks can be nonzero.
\end{problem}
\begin{solution}
Choose independent column lists $a_1,\ldots,a_r$ from $A$ and $d_1,\ldots,d_s$ from $D$. In the first whole matrix, choose the columns $(a_i,0)^\top$ and the columns $(b_j,d_j)^\top$ with the selected lower blocks. In a zero combination, the lower component and independence of the $d_j$ force those $s$ coefficients zero. The upper component and independence of the $a_i$ then force the remaining $r$ coefficients zero. Thus at least $r+s$ columns are independent. Transposing proves the second inequality. For a counterexample with no zero block, $\begin{pmatrix}1&1\\1&1\end{pmatrix}$ has rank one although both diagonal blocks have rank one.
\end{solution}

\begin{problem}[Moving a pivot block]\label{ex:05:off-diagonal-rank}
\begin{enumerate}
\item If $B$ is square and invertible, prove
$\operatorname{rank}\begin{pmatrix}A&B\\0&D\end{pmatrix}=\operatorname{rank}B+\operatorname{rank}(DB^{-1}A)$.
If $C$ is square and invertible, prove
$\operatorname{rank}\begin{pmatrix}A&0\\C&D\end{pmatrix}=\operatorname{rank}C+\operatorname{rank}(AC^{-1}D)$.
\item Set $X=\begin{pmatrix}A&B\\C&0\end{pmatrix}$ and $Y=\begin{pmatrix}0&B\\C&D\end{pmatrix}$. Prove that both ranks equal $\operatorname{rank}B+\operatorname{rank}C$ if either $B$ or $C$ is square and invertible, and that both ranks are at least that sum without this assumption.
\item Prove $\operatorname{rank}X=\operatorname{rank}A+\operatorname{rank}(CA^{-1}B)$ when $A$ is invertible, and $\operatorname{rank}Y=\operatorname{rank}D+\operatorname{rank}(BD^{-1}C)$ when $D$ is invertible.
\end{enumerate}
\end{problem}
\begin{solution}
(a) Swap the two column groups in the first matrix to obtain $\begin{pmatrix}B&A\\D&0\end{pmatrix}$. Its Schur complement at the invertible block $B$ is $-DB^{-1}A$, whose rank equals that of $DB^{-1}A$. For the second matrix, swap the row groups to obtain $\begin{pmatrix}C&D\\A&0\end{pmatrix}$ and use the same argument with Schur complement $-AC^{-1}D$.
(b) Swapping column groups turns $X$ into $\begin{pmatrix}B&A\\0&C\end{pmatrix}$ and $Y$ into $\begin{pmatrix}B&0\\D&C\end{pmatrix}$. The zero-block rank equality under an invertible diagonal block and the preceding lower bounds give exactly the assertions.
(c) Eliminate using $A$ in $X$ or $D$ in $Y$. The remaining Schur blocks are $-CA^{-1}B$ and $-BD^{-1}C$, respectively; multiplication by $-1$ does not change rank.
\end{solution}

\begin{problem}[Frobenius and Sylvester inequalities]\label{ex:05:rank-inequalities}
\begin{enumerate}
\item For a defined triple product $ABC$, prove
$\operatorname{rank}(AB)+\operatorname{rank}(BC)\leqslant\operatorname{rank}B+\operatorname{rank}(ABC)$ using block matrices.
\item For $A:m\times p$, $B:p\times n$, deduce $\operatorname{rank}(AB)\geqslant\operatorname{rank}A+\operatorname{rank}B-p$.
\item Deduce that $AB=0$ implies $\operatorname{rank}A\leqslant p-\operatorname{rank}B$.
\end{enumerate}
\end{problem}
\begin{solution}
(a) Write $A:m\times n$, $B:n\times p$, $C:p\times q$ and set $Z=\begin{pmatrix}0&AB\\BC&B\end{pmatrix}$. The block identity
\[
 \begin{pmatrix}I_m&-A\\0&I_n\end{pmatrix}
 Z
 \begin{pmatrix}I_q&0\\-C&I_p\end{pmatrix}
 =\begin{pmatrix}-ABC&0\\0&B\end{pmatrix}
\]
is verified by two multiplications. Its multipliers are invertible, so $\operatorname{rank}Z=\operatorname{rank}(ABC)+\operatorname{rank}B$. The lower bound for a zero diagonal block gives $\operatorname{rank}Z\geqslant\operatorname{rank}(AB)+\operatorname{rank}(BC)$.
(b) Apply (a) to the triple $(A,I_p,B)$ and rearrange.
(c) Substitute $\operatorname{rank}(AB)=0$ into (b).
\end{solution}

\begin{problem}[Two rank formulas for one block matrix]\label{ex:05:schur-ranks}
\begin{enumerate}
\item If $A$ is invertible, prove
\[ \operatorname{rank}Z=\operatorname{rank}A+\operatorname{rank}(D-CA^{-1}B). \]
Prove the analogous formula pivoting on invertible $D$.
\item For $B:m\times n$, prove
\[
 \operatorname{rank}\begin{pmatrix}I_m&B\\B^\top&I_n\end{pmatrix}
 =m+\operatorname{rank}(I_n-B^\top B)
 =n+\operatorname{rank}(I_m-BB^\top).
\]
\item For $B:m\times n$, $C:n\times m$, prove
$\operatorname{rank}(I_m-BC)=\operatorname{rank}(I_n-CB)+m-n$.
Include the square case explicitly.
\end{enumerate}
\end{problem}
\begin{solution}
(a) The two Schur eliminations use invertible multipliers and give diagonal block matrices, so rank is the sum of the indicated block ranks.
(b) Apply (a) to the two identity diagonal blocks of the displayed matrix.
(c) Similarly, for $Z=\begin{pmatrix}I_m&B\\C&I_n\end{pmatrix}$, pivoting on $I_m$ gives rank $m+\operatorname{rank}(I_n-CB)$, while pivoting on $I_n$ gives $n+\operatorname{rank}(I_m-BC)$. Equating and rearranging gives the formula. If $m=n$, it becomes equality of the two ranks. The whole block matrix has an additional identity-block rank; that term must not be omitted.
\end{solution}

\begin{problem}[Another proof of the sum-rank bound]\label{ex:05:block-sum-rank}
Use partitioned matrices to prove $\operatorname{rank}(A+B)\leqslant\operatorname{rank}A+\operatorname{rank}B$ for $A,B\in\R^{m\times n}$.
\end{problem}
\begin{solution}
The identity
\[
 A+B=(I_m\ I_m)\begin{pmatrix}A&0\\0&B\end{pmatrix}\begin{pmatrix}I_n\\I_n\end{pmatrix}
\]
expresses the sum as a product through a block diagonal matrix. Product rank cannot exceed the rank of that middle factor, which is $\operatorname{rank}A+\operatorname{rank}B$. This is a block proof with no assumption that either matrix is square.
\end{solution}

\subsection*{Optional computational practice: sweeping}
For a square matrix $A=(a_{ij})$ and an index $k$ with $a_{kk}\ne0$, define $B=\operatorname{SWP}(A,k)$ by
\[
 b_{kk}=-\frac1{a_{kk}},\quad b_{ik}=\frac{a_{ik}}{a_{kk}},\quad
 b_{kj}=\frac{a_{kj}}{a_{kk}},\quad
 b_{ij}=a_{ij}-\frac{a_{ik}a_{kj}}{a_{kk}}
 \quad(i,j\ne k).
\]
All entries on the right refer to the matrix before that sweep.

\begin{problem}[Two successive sweeps]\label{ex:05:sweep-two}
For $A=\begin{pmatrix}a&b\\c&d\end{pmatrix}$, compute $A^{(1)}=\operatorname{SWP}(A,1)$ and $A^{(2)}=\operatorname{SWP}(A^{(1)},2)$, state all nonzero-pivot conditions, and prove $A^{(2)}=-A^{-1}$ when both sweeps are defined.
\end{problem}
\begin{solution}
The first sweep requires $a\ne0$ and gives
\[
 A^{(1)}=\begin{pmatrix}-1/a&b/a\\c/a&d-bc/a\end{pmatrix}.
\]
The second additionally requires $d-bc/a\ne0$, equivalently $ad-bc\ne0$. Substitution in the definition gives
\[
 A^{(2)}=-\frac1{ad-bc}\begin{pmatrix}d&-b\\-c&a\end{pmatrix}.
\]
For example, its upper-left entry is $-1/a-(bc/a^2)/(d-bc/a)=-d/(ad-bc)$; the other three entries follow directly from the pivot row and column rules. Multiplying the final displayed matrix by $A$ in either order gives $-I_2$, proving the inverse claim without a determinant theorem.
\end{solution}

\begin{problem}[Sweeping an interior index]\label{ex:05:sweep-general}
Compute $\operatorname{SWP}(A,2)$ for a general $3\times3$ matrix and give the corresponding block formula for $\operatorname{SWP}(A,p)$ in an $n\times n$ matrix. State the condition for each formula.
\end{problem}
\begin{solution}
For $3\times3$, the answer is
\[
 \begin{pmatrix}
 a_{11}-a_{12}a_{21}/a_{22}&a_{12}/a_{22}&a_{13}-a_{12}a_{23}/a_{22}\\
 a_{21}/a_{22}&-1/a_{22}&a_{23}/a_{22}\\
 a_{31}-a_{32}a_{21}/a_{22}&a_{32}/a_{22}&a_{33}-a_{32}a_{23}/a_{22}
 \end{pmatrix},
\]
requiring $a_{22}\ne0$. For the general case, group indices before $p$, the single index $p$, and indices after $p$, and write
\[
 A=\begin{pmatrix}A_{11}&u&A_{13}\\v^\top&a&w^\top\\A_{31}&z&A_{33}\end{pmatrix}.
\]
Here $a=a_{pp}$; the result, for $a\ne0$, is
\[
 \begin{pmatrix}
 A_{11}-uv^\top/a&u/a&A_{13}-uw^\top/a\\
 v^\top/a&-1/a&w^\top/a\\
 A_{31}-zv^\top/a&z/a&A_{33}-zw^\top/a
 \end{pmatrix}.
\]
Each block is obtained by collecting exactly the entry updates in the definition. For $p=1$ or $p=n$, the absent before/after block groups are simply omitted.
\end{solution}

\begin{problem}[The sweeping theorem]\label{ex:05:sweep-theorem}
Starting with $A^{(0)}=A\in\R^{n\times n}$, define $A^{(k)}=\operatorname{SWP}(A^{(k-1)},k)$ for $k=1,\ldots,p$, assuming every required pivot is nonzero. With the original partition $A=\begin{pmatrix}P&Q\\R&S\end{pmatrix}$ and $P:p\times p$, prove
\[
 A^{(p)}=\begin{pmatrix}-P^{-1}&P^{-1}Q\\RP^{-1}&S-RP^{-1}Q\end{pmatrix}.
\]
Deduce that a complete defined sequence of sweeps gives $A^{(n)}=-A^{-1}$.
\end{problem}
\begin{solution}
View $y=Ax$ as a relation between coordinates. Solving its $k$th equation for $x_k$ shows that replacing the input coordinate $x_k$ by $y_k$ and the output coordinate $y_k$ by $-x_k$ changes its matrix exactly according to the sweep definition. Indeed,
\[
 -x_k=-y_k/a_{kk}+\sum_{j\ne k}(a_{kj}/a_{kk})x_j,
\]
and substituting $x_k$ in each other equation gives coefficient $a_{ik}/a_{kk}$ on $y_k$ and coefficient $a_{ij}-a_{ik}a_{kj}/a_{kk}$ on $x_j$.
After the first $p$ sweeps, the new input is $(y_1,\ldots,y_p,x_{p+1},\ldots,x_n)$ and the new output is $(-x_1,\ldots,-x_p,y_{p+1},\ldots,y_n)$. The hypotheses ensure this is a well-defined linear rule for every new input. The equation $Px=0$ therefore has only the zero solution: set the unswept input coordinates to zero and the first $p$ original output coordinates to zero, so the new input is zero and hence its new output, including $-x$, is zero. Thus $P$ is invertible.
Writing the original input and output in blocks as $x=(u,v)$ and $y=(r,s)$, the original equations now give
\[
 -u=-P^{-1}r+P^{-1}Qv,\qquad
 s=RP^{-1}r+(S-RP^{-1}Q)v.
\]
These are exactly the asserted blocks. When $p=n$, there are no unswept coordinates and the matrix is $-A^{-1}$.
\end{solution}

\begin{problem}[Solving equations by sweeping]\label{ex:05:sweep-system}
Explain how a defined sequence of sweeps can solve $PX=Q$ for an invertible square $P$. Explain how equation-row permutations can arrange usable pivots if the original fixed order fails. Solve $2x_1+3x_2=8$, $4x_1+5x_2=14$ using sweeps.
\end{problem}
\begin{solution}
Place $P,Q$ in the upper block row of a square matrix $\begin{pmatrix}P&Q\\0&0\end{pmatrix}$, with lower block sizes chosen to make it square. Sweeping the indices of $P$ gives upper-right block $P^{-1}Q$ by the theorem. A prescribed fixed sweep order may fail even for invertible $P$; an equation-row permutation must be applied to both $P$ and $Q$ to arrange usable successive elimination pivots. Ordinary Gaussian pivot selection supplies such a permutation for invertible $P$.
For the numerical system, start with
\[
 A^{(0)}=\begin{pmatrix}2&3&8\\4&5&14\\0&0&0\end{pmatrix}.
\]
The first pivot is $2$, and the sweep gives
\[
 A^{(1)}=\begin{pmatrix}-1/2&3/2&4\\2&-1&-2\\0&0&0\end{pmatrix}.
\]
The second pivot is $-1$, and the next sweep gives
\[
 A^{(2)}=\begin{pmatrix}5/2&-3/2&1\\-2&1&2\\0&0&0\end{pmatrix}.
\]
The upper-right column is $(1,2)^\top$, the solution. Substitution gives $2+6=8$ and $4+10=14$.
\end{solution}
